\section{Investigaciones previas y brecha de conocimiento}
Estos antecedentes muestran que las rúbricas pueden utilizarse tanto para evaluar respuestas abiertas como para orientar su generación durante el entrenamiento. En esta línea, Rubicon estudia su uso como señal de recompensa~\cite{huang2025rubicon}, mientras que ARES automatiza la construcción de preguntas, respuestas y criterios a partir de documentos~\cite{li2026ares} y RuscaRL incorpora ayudas basadas en rúbricas que se retiran progresivamente~\cite{zhou2026ruscarl}. A partir de estos enfoques, este trabajo propone una metodología que aprovecha pares de consulta y respuesta existentes para realizar una adaptación supervisada inicial y extraer, de una partición reservada, los criterios del entrenamiento posterior. Para integrar esos criterios en la optimización, la propuesta combina comprobaciones mediante código con un modelo de lenguaje que actúa como juez semántico para evaluar los conceptos no detectados por esas comprobaciones, junto con la retirada progresiva de las entidades mostradas como ayuda durante GRPO. De esta manera, el aporte se concentra en combinar y adaptar estos mecanismos dentro de un procedimiento de adaptación de dominio, estudiado en consultas clínicas abiertas.

\subsection{Rúbricas como señales de recompensa}
Una línea de investigación reciente busca extender el aprendizaje por refuerzo a tareas abiertas mediante recompensas basadas en rúbricas, que descomponen la calidad de una respuesta en criterios explícitos. Dentro de esta línea, Huang y sus colaboradores presentaron Rubicon, un enfoque que utiliza rúbricas elaboradas por personas, generadas por modelos de lenguaje o construidas mediante una combinación de ambos para orientar el entrenamiento en tareas de escritura, humanidades y seguimiento de instrucciones. Su propuesta combina criterios generales con otros específicos de cada tarea o consulta, permitiendo evaluar propiedades que no pueden reducirse a la comprobación de una única respuesta correcta. Para ello, construyeron un banco de más de 10.000 rúbricas y entrenaron Qwen3-30B-A3B con algo más de 5.000 ejemplos, obteniendo un aumento del puntaje promedio de 65,29 a 70,50 en las siete evaluaciones abiertas reportadas, aunque el puntaje en IFEval descendió de 83,55 a 81,70 \cite{huang2025rubicon}. Estos resultados muestran mejoras en el promedio junto con retrocesos en una evaluación particular, y los autores señalan que incorporar más rúbricas no garantiza por sí solo mejores resultados, ya que también importan su diversidad, su nivel de detalle y la estrategia utilizada durante el entrenamiento. Este antecedente resulta relevante para el presente trabajo porque muestra cómo las rúbricas pueden emplearse como señales de optimización, aunque su aplicación a consultas clínicas requiere considerar criterios y riesgos propios del dominio.

%
\subsection{Generación automática de rúbricas}
En relación con la construcción automática de estos criterios, Li y sus colaboradores propusieron ARES, un método que transforma documentos en preguntas, respuestas de referencia y rúbricas específicas para cada ejemplo, con el objetivo de generar datos para el entrenamiento por refuerzo en tareas abiertas \cite{li2026ares}. Para ello, el proceso incorpora distintas perspectivas de usuario y características del dominio durante la generación de las preguntas, además de controles para comprobar que puedan responderse sin acceder al documento original, que las respuestas estén respaldadas por su contenido y que las rúbricas cubran los aspectos relevantes de cada consulta. A partir de este procedimiento, los autores construyeron un conjunto de 101.847 ejemplos y 1.108.163 criterios distribuidos en diez dominios, que utilizaron para estudiar el entrenamiento de Qwen3-4B y comparar distintas señales de recompensa, entre ellas las rúbricas específicas, las rúbricas generales y las respuestas de referencia. En cuanto a los resultados, la variante ARES-RL obtuvo el mejor promedio en las siete evaluaciones consideradas, aunque esa ventaja no se mantuvo en todos los escenarios, ya que en la comparación de recompensas sobre HealthBench alcanzó un puntaje de 41,45, frente a 45,97 con rúbricas generales y 50,13 con respuestas de referencia \cite{li2026ares}. El artículo identifica esta evaluación como HealthBench, pero no especifica si corresponde al conjunto original, a HealthBench Hard o a HealthBench Consensus, y señala que las tareas abiertas se evalúan mediante los protocolos de rúbricas o jueces propios de cada benchmark, sin detallar el modelo evaluador de HealthBench ni la fórmula de agregación utilizada para obtener su puntaje global. Por ello, no puede comprobarse a partir de esa descripción si reproduce íntegramente el cálculo oficial. De esta forma, la investigación de ARES muestra cómo generar rúbricas a partir de documentos, pero sus resultados también indican que los criterios específicos de cada consulta no siempre ofrecen la mejor señal de entrenamiento.

%
\subsection{Rúbricas como ayudas durante el entrenamiento}
Además de construir la recompensa, otra cuestión es cómo lograr que el modelo genere respuestas suficientemente buenas para aprender de ellas durante el entrenamiento. Zhou y sus colaboradores abordaron este problema mediante RuscaRL, un método que utiliza las rúbricas tanto para orientar la generación como para evaluar las respuestas \cite{zhou2026ruscarl}. Para ello, proporcionan distintos niveles de ayuda dentro de un mismo grupo de generaciones y reducen progresivamente la información mostrada al modelo a medida que avanza el entrenamiento, buscando favorecer la exploración inicial y disminuir luego la dependencia de esas ayudas. Esta retirada afecta a los criterios incluidos en la consulta, mientras que la rúbrica completa se conserva para calcular la recompensa mediante un modelo evaluador que comprueba el cumplimiento de cada criterio. En sus experimentos con Qwen2.5-7B-Instruct, los autores evaluaron sobre HealthBench-500, un subconjunto de 500 casos seleccionados aleatoriamente de HealthBench, y utilizaron los 4.500 restantes para entrenamiento, una partición que permite identificar el conjunto original del benchmark, aunque el artículo no lo denomina expresamente así frente a HealthBench Hard y HealthBench Consensus. Bajo ese protocolo, el puntaje pasó de 23,4 para el modelo inicial a 56,4 con RuscaRL, frente a 52,0 con entrenamiento por refuerzo basado en rúbricas sin el mecanismo de ayudas. Para esta evaluación, los autores utilizaron GPT-4.1 como juez, con el formato de evaluación por criterio de HealthBench, y reportaron el promedio de tres ejecuciones en escala porcentual \cite{zhou2026ruscarl}. Aunque describen la suma ponderada y la normalización de los criterios para la recompensa de entrenamiento, no explicitan por separado la fórmula completa del puntaje global del benchmark ni el tratamiento de los valores negativos, por lo que esa descripción no permite confirmar que reproduzcan íntegramente el cálculo oficial. Esta comparación permite examinar el aporte del andamiaje más allá del uso de rúbricas como recompensa, aunque debe considerarse que el entrenamiento médico también utilizó casos de HealthBench distintos de los reservados para evaluación. Así, el trabajo ofrece un antecedente para retirar progresivamente las ayudas, pero queda abierta su aplicación cuando los datos de entrenamiento provienen de un corpus clínico independiente del benchmark.

%
